\section{特征学习与核方法的差异}
\label{sec:14-Deep-Learning/07-Interpretability-Mathematics/02}

核方法和深度网络都可以理解为先选择表示，再在表示上完成预测。关键差异不在``一个线性、一个非线性''：带通用核的方法同样能逼近很丰富的函数。更有用的区别是，标准核方法在拟合预测器时通常保持特征几何不变，而深度网络会让表示随任务损失共同变化。

\begin{quote}
\textbf{先抓住这一点：}特征学习改变的是``哪些输入彼此接近、哪些方向容易分开''。它带来适应任务的可能，也把表示偏差、优化困难和数据捷径一起带进模型。
\end{quote}

\subsection{固定几何与可训练几何}

核函数 \(K(x,x')\) 定义输入之间的相似性；若存在特征映射 \(\phi\)，则

\[
K(x,x')=\ip{\phi(x)}{\phi(x')}.
\]

标准核回归或支持向量机在这个几何中学习预测器，训练不会因为当前标签而任意改写 \(\phi\)。核及其超参数当然可以由数据选择，深核和多核学习也会学习几何；因此``核永远固定、从不利用数据''并不准确。这里比较的是固定核制度与端到端可训练表示，而不是给所有核方法下结论。

深度网络把表示写成 \(\phi_\theta(x)\)，并与最后的预测头 \(w\) 一起优化。更新 \(\theta\) 会改变样本间距离、可分方向和局部不变性。卷积预先写入局部参数共享，注意力预先写入内容相关读取；训练再从这些允许的结构中选择适合当前目标的表示。

\subsection{能力差异不等于表示定理差异}

固定的通用核在适当条件下也能逼近广泛函数，因此不能用``核方法表示不了、深度网络表示得了''概括实践差异。更常见的区别来自：

\begin{itemize}
\tightlist
\item 哪种函数在当前表示中复杂度较低；
\item 样本量增加时需要付出多少计算和存储；
\item 结构先验是否与图像、序列或图等数据形式匹配；
\item 优化过程是否显著改变中间表示；
\item 学到的相似性在分布变化后是否仍有意义。
\end{itemize}

比较方法时应控制参数量、训练预算、数据增强和核选择。单次准确率差异不能单独证明``发生了特征学习''。

\subsection{怎样观察特征是否真的在学习}

可以比较初始化与训练后表示的相似度，冻结不同层做线性探测，或者把完整训练与只训练最后一层、对应 NTK 基线作对照。还应通过输入扰动和跨域迁移检查表示利用的是预期结构还是训练捷径。

这些诊断仍不是因果证明。线性探测表现提高，可能来自信息被重新排列，也可能只是维度和正则化改变。可靠结论需要明确对照组和干预对象。

\subsection{低维结构只是候选条件}

若观测由少量潜在因素生成，或任务只依赖少数变化方向，学习表示可能把这些方向变得更容易读取。但低内在维度既不是深度网络成功的充分条件，也不是必要条件：低维流形可能具有复杂决策边界，高维问题也可能因对称性、局部性或组合结构而可学。

网络学到低维表示也不能自动说明它找到了``真实因素''。监督目标可能压缩掉与标签无关但现实上重要的信息，或者保留一个只在训练环境中有效的捷径。

\subsection{什么时候固定核更合适}

当样本较少、核结构已有可靠领域依据、需要凸优化或清晰的不确定性表达时，固定核可能更省数据、更稳定，也更容易审计。端到端特征学习在数据和计算充足、结构先验能写入架构且迁移评价充分时更有吸引力。

真正的选择不是``传统方法还是深度方法''，而是：相似性应当预先规定多少，又有多少证据允许它由当前数据学习。下一篇进一步讨论怎样检验所谓低维结构，而不把它当作训练成功后的事后解释。
